% Section 9.2, from lectures/L09/appendix/b_loops_and_delays.md.

\section{Loopar och tidsfördröjningar}
\label{c9:app:b}

\subsection{En loop med ett bestämt antal varv}
\label{c9:b1}
En \textbf{loop} är en del av programmet som körs om och om igen. Den byggs med ett hopp bakåt,
till en etikett före det som ska upprepas. Oftast ska loopen gå ett bestämt antal varv, och då
behövs en \textbf{räknare}: ett register som räknas ned för varje varv, tills det blir noll.

Programmet \code{count_loop.asm} räknar ut summan $1 + 2 + \dots + 10$:

\bookfigure{l09_count_loop}{Flödesplan för en räknarloop.}{c9:fig:count_loop}

\begin{asmcode}
    clr r17                         ; r17 = the sum so far = 0
    ldi r16, 10                     ; r16 = the counter: laps left to run

loop:
    out PORTB, r16                  ; Show the counter.
    add r17, r16                    ; sum = sum + counter
    dec r16                         ; One lap fewer left. Sets Z when it reaches 0.
    brne loop                       ; Not zero yet: go round again.
\end{asmcode}

\noindent Nyckeln är de två sista raderna. \code{dec r16} minskar räknaren och sätter \code{Z = 1}
när den når noll. \code{brne loop} hoppar tillbaka så länge \code{Z = 0}, alltså så länge räknaren
inte är noll. Ingen \code{cpi} behövs: \code{dec} sätter själv den flagga hoppet läser.

Följ räknaren: den börjar på 10, och \code{brne} hoppar tillbaka när den är 9, 8, \dots, 1. När
\code{dec} gör den till 0 hoppar \code{brne} inte, och programmet fortsätter efter loopen. Loopen
går alltså \textbf{10 varv}, lika många som räknarens startvärde, och summan blir
$10 + 9 + \dots + 1 = 55$.

\textbf{En fälla: startvärdet 0.} Om räknaren börjar på 0 gör den första \code{dec} den till 255,
och loopen går \textbf{256 varv}, inte noll. En loop med testet sist körs alltid minst en gång.

\subsection{Testet först eller sist}
\label{c9:b2}
Loopen i \secref{c9:b1} har testet \textbf{sist}: först görs arbetet, sedan frågar \code{brne} om
det ska göras igen. Det är den kortaste formen, och den som används mest i assembler.

Om loopen ska kunna gå \textbf{noll} varv måste testet stå \textbf{först}:

\begin{asmcode}
    tst r16                         ; Is the counter already zero?
loop:
    breq done                       ; Yes: leave before doing any work.
    add r17, r16
    dec r16                         ; Sets Z for the breq at the top.
    rjmp loop
done:
\end{asmcode}

\noindent Nu hoppar \code{breq} ut direkt om räknaren är noll från början. Priset är en instruktion
till per varv, \code{rjmp}, och en klockcykel till: \code{breq} som inte hoppar och \code{rjmp} tar
1 + 2 cykler, mot 2 för \code{brne} som hoppar. Välj formen efter om noll varv kan förekomma.

\subsection{Vad en instruktion kostar i tid}
\label{c9:b3}
Ett program som tänder en lysdiod och sedan släcker den gör det på en bråkdel av en mikrosekund.
För att ögat ska hinna se något, eller för att ett relä ska hinna dra, eller för att ett trafikljus
ska lysa grönt i några sekunder, måste programmet \textbf{vänta}. En \textbf{tidsfördröjning} är
en loop som inte gör något annat än att ta tid.

Hur mycket tid den tar går att räkna ut exakt, eftersom varje instruktion tar ett fast antal
klockcykler (\secref{c7:a7}):

\begin{booktable}{@{}p{6.4em}L@{}}
  \thead{Klockcykler} & \thead{Instruktioner} \\
  \midrule
  1 & \code{ldi} \code{mov} \code{add} \code{sub} \code{subi} \code{inc} \code{dec} \code{and}
    \code{or} \code{eor} \code{com} \code{cp} \code{cpi} \code{lsl} \code{lsr} \code{rol}
    \code{ror} \code{in} \code{out} \code{nop} \\
  2 & \code{rjmp} \code{sbi} \code{cbi} \\
  1 eller 2 & \code{breq} \code{brne} \code{brlo} \code{brsh} och de andra villkorliga hoppen: 1
    om de inte hoppar, 2 om de hoppar \\
\end{booktable}

\noindent Vid 16~MHz tar en klockcykel 62,5~ns, och 16 cykler tar en mikrosekund. Hela tabellen, med
fler instruktioner, finns på referensbladet i bilaga~\ref{app:avr}.

\subsection{En kort tidsfördröjning}
\label{c9:b4}
Den enklaste fördröjningen är en loop som räknar ned ett register, som i \code{short_delay.asm}:

\begin{asmcode}
delay_start:
    ldi r18, 200                    ; 1 cycle
delay_loop:
    dec r18                         ; 1 cycle per lap
    brne delay_loop                 ; 2 cycles per lap, 1 on the last
delay_end:
\end{asmcode}

\noindent Räkna ut kostnaden med en rad per instruktion: vad den kostar, hur många gånger den körs,
och produkten. \code{brne} får två rader, eftersom den kostar olika mycket när den hoppar och när
den inte gör det.

\begin{narrowtable}{@{}lccc@{}}
  \thead{Instruktion} & \thead{Cykler} & \thead{Gånger} & \thead{Summa} \\
  \midrule
  \code{ldi r18, 200} & 1 & 1 & 1 \\
  \code{dec r18} & 1 & 200 & 200 \\
  \code{brne}, hoppar & 2 & 199 & 398 \\
  \code{brne}, hoppar inte & 1 & 1 & 1 \\
  \midrule
  \textbf{Totalt} & & & \textbf{600} \\
\end{narrowtable}

\noindent 600 cykler är $600 \cdot 62,5$~ns = \textbf{37,5~µs}. Det svåra i tabellen är
kolumnen \emph{Gånger}: \code{brne} körs 200 gånger, men hoppar bara 199 av dem, eftersom den
sista gången, när räknaren har blivit noll, går programmet vidare.

Med ett annat antal varv, $n$, blir summan $1 + n + 2(n - 1) + 1 = \mathbf{3n}$ klockcykler. Det
är en formel värd att komma ihåg:
\begin{itemize}
  \item \textbf{Önskad tid $\rightarrow$ antal varv.} 30~µs är $30 \cdot 16 = 480$ cykler, och
    $n = 480 / 3 = 160$.
  \item \textbf{Längsta möjliga.} Ett register rymmer högst 255, men startvärdet 0 ger 256 varv
    (\secref{c9:b1}): 768 cykler, 48~µs. Längre än så räcker inte en loop.
\end{itemize}

\subsection{En längre tidsfördröjning}
\label{c9:b5}
För att vänta längre läggs en loop \textbf{inuti} en annan. Den inre loopen är fördröjningen från
\secref{c9:b4}. Den yttre upprepar den ett antal gånger:

\bookfigure{l09_nested_delay}{Flödesplan för två loopar i varandra.}{c9:fig:nested_delay}

\noindent Programmet \code{long_delay.asm}:

\begin{asmcode}
delay_start:
    ldi r19, OUTER                  ; 1 cycle
delay_outer:
    ldi r18, INNER                  ; 1 cycle per outer lap
delay_inner:
    dec r18                         ; 1 cycle per inner lap
    brne delay_inner                ; 2 per inner lap, 1 on the last
    dec r19                         ; 1 cycle per outer lap
    brne delay_outer                ; 2 per outer lap, 1 on the last
delay_end:
\end{asmcode}

\noindent Räkna ut det i två steg:
\begin{enumerate}
  \item \textbf{Ett varv i den yttre loopen.} \code{ldi r18} (1), den inre loopen utan sin
    \code{ldi} ($3 \cdot \text{INNER} - 1$), \code{dec r19} (1) och \code{brne} som hoppar (2):
    $3 \cdot \text{INNER} + 3$ cykler. Det sista yttre varvet är en cykel kortare, eftersom
    \code{brne} då inte hoppar.
  \item \textbf{Hela fördröjningen.} \code{ldi r19} (1), plus OUTER varv, minus den cykel det
    sista varvet sparar:
\end{enumerate}
\[
  \text{cykler} = \text{OUTER} \cdot (3 \cdot \text{INNER} + 3)
\]

\noindent Med OUTER = 208 och INNER = 255 blir det $208 \cdot 768 = \mathbf{159\,744}$
\textbf{cykler}, vilket är $159\,744 \cdot 62,5$~ns = \textbf{9,984~ms}. Nästan 10~ms.

\subsubsection{Att träffa exakt}
Formeln ger bara vissa tider. 10~ms är 160\,000 cykler, och 160\,000 är inte delbart med 3, så
ingen kombination av OUTER och INNER ger exakt 10~ms. Lösningen är att lägga till en \code{nop},
som inte gör något annat än att ta en cykel, i den yttre loopen. Då kostar varje yttre varv en
cykel mer:
\[
  \text{cykler} = \text{OUTER} \cdot (3 \cdot \text{INNER} + 4)
\]
och OUTER = 250, INNER = 212 ger $250 \cdot 640 = \mathbf{160\,000}$ \textbf{cykler, exakt 10~ms}.
Det gör du i Labb~3 \hyperref[lab3:d16]{del~16}.

Så väljer du konstanterna för en önskad tid:
\begin{enumerate}
  \item Räkna om tiden till cykler: tid i µs $\cdot$ 16.
  \item Välj INNER så stort som möjligt, 255, och räkna ut OUTER = cykler / 768. Avrunda.
  \item Räkna ut vad det faktiskt blir. Behövs mer precision, prova andra kombinationer, eller
    lägg till \code{nop}.
\end{enumerate}

Med två loopar räcker det till ungefär $256 \cdot 771$ cykler, knappt 12,4~ms. För längre tider, en
halv sekund för en blinkande lysdiod, behövs en loop till eller en räknare med 16 bitar, som i
kapitel~\ref{c10:ch}.

\subsection{Mät med stoppuret}
\label{c9:b6}
Ett räknat värde är en förutsägelse. Kontrollera det med simulatorns stoppur, enligt
\secref{studio:4-4}:
\begin{enumerate}
  \item Kontrollera att \textbf{Frequency} i Processor Status står på \textbf{16~MHz}.
  \item Sätt en brytpunkt på raden \code{delay_start} och en på raden \code{delay_end}.
  \item Kör till den första med \textbf{F5}, nollställ stoppuret, och kör till den andra.
  \item Läs av \textbf{Cycle Counter} eller \textbf{Stop Watch}.
\end{enumerate}

För \code{short_delay.asm} ska stoppuret visa exakt 600 cykler, 37,5~µs, och för
\code{long_delay.asm} 159\,744 cykler. Instruktionen där en brytpunkt sitter har ännu inte körts
när programmet stannar, så mätningen från \code{delay_start} till \code{delay_end} omfattar precis
fördröjningen: \code{ldi} men inte instruktionen efter.

\textbf{Om mätningen inte stämmer med räkningen} är det nästan alltid räkningen som är fel.
Kontrollera först kolumnen \emph{Gånger} för hoppen, och sedan var brytpunkterna sitter: en
brytpunkt en rad för tidigt eller för sent tar med eller utelämnar en instruktion. Att stämma av en
räkning mot en mätning och förklara skillnaden är det kontrolluppgiften, övning~\ref{c9:ex:9},
handlar om.
